% Generated draft: requires mathematical and source review.
\documentclass[11pt]{article}
\usepackage[margin=25mm]{geometry}
\usepackage{fontspec}
\setmainfont{Latin Modern Roman}
\usepackage{amsmath,amssymb,mathtools,unicode-math}
\setmathfont{Latin Modern Math}
\usepackage{enumitem,needspace,xurl,hyperref}
\setlength{\parindent}{0pt}
\setlength{\parskip}{5pt}
\setlength{\emergencystretch}{3em}
\newcommand{\N}{\mathbb{N}}
\newcommand{\Z}{\mathbb{Z}}
\newcommand{\Q}{\mathbb{Q}}
\newcommand{\R}{\mathbb{R}}
\newcommand{\C}{\mathbb{C}}
\newcommand{\F}{\mathbb{F}}
\newcommand{\A}{\mathbb{A}}
\newcommand{\PP}{\mathbb{P}}
\newcommand{\E}{\mathbb{E}}
\newcommand{\eps}{\varepsilon}
\newcommand{\dd}{\,\mathrm{d}}
\begin{document}
\section*{UnsolvedMath ID: 801016}
\textbf{Source problem code:} catalogue entry
\subsection{Definitions and mathematical statement}
The source leaves the domain of \(x\) and the properties of \(U\) unspecified. Under the conventional interpretation of relativized prefix-free Kolmogorov complexity, let \(2^{<\omega}\) be the set of finite binary strings, let \(2^\omega\) be the set of infinite binary sequences, and regard \(x\in 2^\omega\) as an oracle subset of \(\N\). Fix an optimal universal oracle machine \(U\) whose domain is prefix-free for each oracle. For \(\sigma\in 2^{<\omega}\), define
\[
K^x(\sigma)=\min\{\lvert\tau\rvert:\tau\in 2^{<\omega},\ U^x(\tau)=\sigma\},
\qquad
\Omega^x=\sum_{\sigma\in 2^{<\omega}}2^{-K^x(\sigma)}.
\]
Here \(\lvert\tau\rvert\) denotes the length of the program \(\tau\). The prefix-free convention makes this a convergent real-valued sum. The displayed definition sums one term for each output string, using its shortest program length; it is to be retained exactly as supplied.

For sets \(A,B\subseteq\N\), write \(A\leq_T B\) when membership in \(A\) is computable using oracle \(B\). When a real serves as an oracle, use its binary expansion, choosing the terminating expansion at dyadic reals. Let \(0''\) denote the second Turing jump of the empty set.

With these conventions, the question attributed in the supplied record to Downey--Hirschfeldt--Miller--Nies [DHMN] is:
\[
\exists x\in 2^\omega\quad 0''\leq_T\Omega^x\,?
\]
The oracle domain and machine conventions above require confirmation against the underlying source. The supplied record does not specify whether the intended assertion concerns a particular fixed machine or every optimal universal prefix-free oracle machine; no machine-independence assertion is made here.
\subsection{Short English statement}
Under the stated conventional interpretation, does some binary-sequence oracle \(x\) make the real \(\Omega^x\), defined by summing \(2^{-K^x(\sigma)}\) over all finite binary strings, compute the second Turing jump \(0''\)? The source labels this problem open, but that status and the unresolved machine conventions require review.
\subsection{Sources}
\begin{itemize}
\item Supplied local UnsolvedMath registry record, permanent ID \(801016\), source code \texttt{catalogue entry}, in the AMR Open Problem Lists collection. The supplied status is \texttt{open}; this is a registry label, not an independently verified literature assessment.
\item Barmpalias, Gao et al., \emph{Open Problems in Mathematical Logic} (2026), Problem 28. Supplied source URL: \url{https://arxiv.org/abs/2608.26628}. The record reports access on 2026-10-16 and extraction by \texttt{arxiv\_latex}; these are supplied provenance details, not a claim of independent source access.
\item Attribution/context supplied in the title and background: Downey--Hirschfeldt--Miller--Nies [DHMN]. No full bibliographic entry for [DHMN] is included in the supplied record. The registry's proposer field lists Barmpalias, Gao et al.
\end{itemize}
\end{document}
